\documentclass[a4paper, dvipsnames, table]{article} 
%pensavo di usare table per i nomi in copertina
\usepackage[utf8]{inputenc}
\usepackage[T1]{fontenc}
\usepackage[italian]{babel}
\usepackage{graphicx}

\begin{document}
\title{Burgheria Padovana}
\date{Anno accademico: 2020/2021}

\newpage  %includere repo github?
\begin{figure}
	\centering
	\includegraphics[width=2cm]{../img/burger-icon.png}%da sistemare qualità
\end{figure}
\maketitle

\begin{table}
	\centering
	\begin{tabular}{c|c c c}
		\textbf{Componenti}
			& Albiero	& Davide      & 1193425 \\
			& Fincato	& Alessandro  &         \\
 			& Panighel	& Cristiano   & 1201284 \\
			& Tossuto	& Matteo      & 1193493 \\
	\end{tabular}
\end{table}

\begin{center}
	\textbf{Indirizzo sito web}: \\
	\textbf{Email referente del gruppo}: davide.albiero@studenti.unipd.it
\end{center}

\begin{table}
	\centering
	\begin{tabular}{c|c c}
		\textbf{Utenti} & \textbf{Nickname} & \textbf{Password} \\
		\hline
		Admin           & admin & admin             \\
		User            & user  & user              \\
	\end{tabular}
\end{table}

\newpage
	\tableofcontents

\newpage
\section{Introduzione}
	Il progetto \emph{Burgheria Padovana} vuole implementare un sito Internet che offra la possibilità di fornire informazioni riguardo il suo punto vendita.\\
Il locale è aperto già da diversi anni e ha finalmente deciso di rinnovarsi creando un proprio sito internet.\\
Il sito dovrà contenere informazioni riguardanti i panini che offre, suddivisi nelle varie categorie: Pollo, Manzo e Speciali. 
Inoltre conterrà gli eventi a cui sarà possibile partecipare, la storia del locale, gli orari di apertura e come contattare il locale.\\
Il sito permette ad un utente privilegiato (admin) di inserire o eliminare gli eventi e di controllare ed eliminare i commenti degli altri utenti. Gli utenti normali (user) saranno visitatori 
con privilegi minimi per poter inserire od eliminare i propri commenti oltre a poter visitare il sito normalmente.\\
Gli \emph{user} dovranno effettuare la registrazione o il login prima di poter usufruire dei privilegi.\\
Inoltre deve essere garantita l'accessibilità in modo che chiunque possa navigare nel sito senza problemi e tranquillamente.\\
Terminato con l'accessibilità, si pone l'attenzione sull'usabilità: rispettando la separazione tra struttura, presentazione e comportamento e rispettando gli standard \emph{W3C} per quanto riguarda \emph{HTML} e \emph{CSS}.\\
Il sito dovrà garantire una navigazione fluida agli utenti evitando disorientamento e, nel caso accadesse, fornirgli supporto per tornare al sito.\\
	
\newpage
\section{Analisi}
	\subsection{Studio dell'utenza finale}
		Il locale Burgheria Padovana offre un prodotto internazionale, comodo, veloce e buono.\\
Pertanto il sito è pensato per rivolgersi ad una categoria di utenti eterogenei, ad esempio: da consumatori abituali a chi vuole provare qualcosa di nuovo, da chi ha bisogno di mangiare qualcosa al volo a chi vuole gustarsi un buon pasto con calma.\\
Queste categorie di utenti, con privilegi minimi, verrà denominata come \emph{utente generico}.
Mentre l'utente con privilegi verrà denominato come \emph{amministratore}.\\ 
Entrambe le tipologie potranno accedere ai loro privilegi autenticandosi tramite form di login.\\
Essendo un utenza finale generica, sarà necessario utilizzare un linguaggio informale, semplice e comprensibile.
Allo stesso modo si andrà a creare un sito di struttura e layout semplici e simili ai modelli a cui l'\emph{utente generico} è abituato.
Si cercherà, quindi, di non rompere le convenzioni esterne e offrendo, indipendentemente dal browser o dal dispositivo utilizzato, una navigazione veloce e intuitiva.\\
	\subsection{Casi d'uso}
		\subsubsection{Utente generico}
			%cosa può fare l'utente
		\subsubsection{Amministratore}
			%cosa può fare di diverso dagli utenti generici

\newpage
\section{Progettazione}%Obiettivi? Design e DB
	\subsection{Obiettivi}
	\subsection{Design}%layout per differenza tra desktop e mobile?
	\subsection{Database}

\newpage
\section{Presentazione}%Per spiegare scelte generali o particolari nel CSS

\newpage
\section{Implementazione}%Per spiegare scelte per HTML, PHP e JS
	\subsection{Linguaggi}
		\subsection{HTML5 e CSS}
		\subsection{SQL}
		\subsection{PHP}
		\subsection{JavaScript}

\newpage
\section{Validazione}%Strumenti per la validazione

\newpage
\section{Fase di Test e Strumenti}%Dove e con cosa testato
	\subsection{XAMPP}
		\subsection{XAMPP}
\emph{XAMPP} è una distribuzione di Apache che permette, tra le varie funzioni, l'utilizzo di \emph{MySQL} e \emph{PHP} di cui abbiamo usufruito.\\
Il loro utilizzo ci ha permesso di realizzare in locale la parte dinamica del sito e la parte di Database, testando le richieste \emph{PHP} e le interrogazioni \emph{MySQL}.
	\subsection*{Silktide - Toolbar}	
	
\newpage
\section{Suddivisione del lavoro}

\end{document}